\subsection{线性映射的转置}

设E、F为两个左A-模；对每个线性映射 \(u: E \to F\) ，映射 \(\operatorname{Hom}(u, l_{A_s})\) 是右A-模 \(F^*\) 到右A-模 \(E^*\) 中的一个线性映射（§1，第2号），称为u的转置。

换言之：

定义5. 对于每个线性映射 \(u\) （从A-模 \(E\) 到A-模 \(F\) 中），线性映射 \(y^{*} \mapsto y^{*} \circ u\) （从 \(F^{*}\) （ \(F\) 的对偶）到 \(E^{*}\) （ \(E\) 的对偶）中）称为 \(u\) 的转置，记作 \(^{t}u\) 。

转置 \(^{t}u\) 因此由关系定义

\[
\langle u (x), y ^ {*} \rangle = \langle x, ^ {t} u (y ^ {*}) \rangle \quad \text { for   all } x \in \mathrm{E} \text { and   all } y ^ {*} \in \mathrm{F} ^ {*}.\tag{15}
\]

定义5不加修改地适用于右A-模，并且此时等价于关系

\[
\langle y ^ {*}, u (x) \rangle = \langle^ {t} u (y ^ {*}), x \rangle \quad \text { for   all } x \in \mathrm{E} \text { and   all } y ^ {*} \in \mathrm{F} ^ {*}.
\]

这里§1第2号中的公式(9)和(10)给出

\[
{ } ^ { t } ( u _ { 1 } + u _ { 2 } ) = { } ^ { t } u _ { 1 } + { } ^ { t } u _ { 2 }\tag{16}
\]

对于两个元素 \(u_{1}, u_{2}\) （属于 \(\operatorname{Hom}_{\mathbf{A}}(\mathbf{E}, \mathbf{F})\) ）以及

\[
{ } ^ { t } ( v \circ u ) = { } ^ { t } u \circ { } ^ { t } v\tag{17}
\]

对于 \(u \in \operatorname{Hom}_{\mathbf{A}}(\mathbf{E}, \mathbf{F})\) 和 \(v \in \operatorname{Hom}_{\mathbf{A}}(\mathbf{F}, \mathbf{G})\) , \(\mathbf{G}\) 作为第三个A模；最后，显然

\[
{ } ^ { t } \mathbf { l } _ { \mathbf { E } } = \mathbf { l } _ { \mathbf { E } ^ { * } } .\tag{18}
\]

注：由(17)和(18)可知，若 \(u\) 是左（相应地，右）可逆的， \(^{t}u\) 是右（相应地，左）可逆的。

命题8. 设 \(u: \mathbf{E} \to \mathbf{F}\) 是一个A-线性映射， \(\mathbf{M}\) 是 \(\mathbf{E}\) 的一个子模， \(\mathbf{M}'\) 是 \(\mathbf{M}\) 在 \(\mathbf{E}^*\) 中的正交；则 \(u(\mathbf{M})\) 在 \(\mathbf{F}^*\) 中的正交是逆像 \(^t u^{-1}(\mathbf{M}')\) .

这由（15）立即得出。

推论. \(u(\mathbf{E})\) 在 \(\mathbf{F}^*\) 中的正交是核 \({}^t u^{-1}(0)\) （即 \({}^t u\) 的核）.

\(\mathbf{E}\) 在 \(\mathbf{E}^*\) 中的正交是0。

若 \(u: \mathbf{E} \to \mathbf{F}\) 是一个同构，则 \({}^{t}u: \mathbf{F}^{*} \to \mathbf{E}^{*}\) 是一个同构，且若 \(v: \mathbf{F} \to \mathbf{E}\) 是 \(u\) 的逆同构，则 \({}^{t}v\) 是 \({}^{t}u\) 的逆同构（公式（17）和（18））。

给定A-模E到A-模F的一个同构u，u的逆同构的转置（等于u的转置的逆同构）称为u的逆步同构，并记为 \(\check{u}\).

该同构 \(\check{u}\) 因此由以下关系所刻画：

\[
\langle u (x), \check {u} (y ^ {*}) \rangle = \langle x, x ^ {*} \rangle \quad \text { for } x \in \mathbf {E}, x ^ {*} \in \mathbf {E} ^ {*}.\tag{19}
\]

若 \(v: \mathbf{F} \to \mathbf{G}\) 是一个同构，则 \(v \circ u\) 的逆步同构是 \(v' \circ \check{u}\) .

特别地，映射 \(u \mapsto \check{u}\) 是线性群 \(\mathbf{GL}(\mathbf{E})\) 到线性群 \(\mathbf{GL}(\mathbf{E}^*)\) 的一个子群上的同构。

设 \(\sigma: \mathbf{A} \to \mathbf{B}\) 是环 \(\mathbf{A}\) 到环 \(\mathbf{B}\) 上的一个同构， \(\mathbf{E}\) 是一个左 \(\mathbf{A}\) -模， \(\mathbf{F}\) 是一个左 \(\mathbf{B}\) -模，而 \(u: \mathbf{E} \to \mathbf{F}\) 是一个关于 \(\sigma\) 的半线性映射（§1，第13号）。设 \(\sigma^{-1}\) 是 \(\sigma\) 的逆同构；对所有 \(y^{*} \in \mathbf{F}^{*}\) ，映射 \(x \mapsto \langle u(x), y^{*} \rangle^{\sigma - 1}\) （从 \(\mathbf{E}\) 到 \(\mathbf{A}\) 中）是一个线性形式；若把它记作 \(^t u(y^*)\) ，就定义了一个映射 \(^t u: \mathbf{F}^{*} \to \mathbf{E}^{*}\) ，它也称为半线性映射 \(u\) 的转置；于是它由恒等式

\[
\langle u (x), y ^ {*} \rangle = \langle x, ^ {t} u (y ^ {*}) \rangle^ {\sigma}\tag{20}
\]

对 \(x \in \mathbf{E}\) , \(y^* \in \mathbf{F}^*\) 成立。立即可以验证， \({}^t u\) 是关于 \(\sigma^{-1}\) 的一个半线性映射。若 \(v\) 表示将 \(u\) 视为 \(\mathbf{E}\) 到 \(\sigma_*(\mathbf{F})\) 中的A-线性映射（§1，第13号）所得的映射，则可写 \(u = \phi \circ v\) ，其中 \(\phi\) 是恒等映射 \(\sigma_*(\mathbf{F}) \to \mathbf{F}\) （视为关于 \(\sigma\) 的半线性映射）。显然 \({}^t u = {}^t v \circ {}^t \phi\) ，且 \(({}^t \phi, \sigma^{-1})\) 是 \(\mathbf{F}^*\) 到 \((\sigma_*(\mathbf{F}))^*\) 上的、关于同构 \(\sigma^{-1}\) 的一个双同构；这一关系使我们能立即把线性映射的转置的性质推广到半线性映射的转置。

